\documentclass[12pt]{report}

% ---------------------------------------------------------------
% 字型與中文排版（和 C++ 演算法解題 1、2 相同）
% ---------------------------------------------------------------
% 作者在 macOS 上用 Times New Roman／Songti TC／PingFang TC；
% 其他系統（例如 Linux）找不到這些字型時，改用 TeX Gyre Termes 與
% Noto Serif／Sans CJK TC，同一份原始碼兩邊都能編。
\usepackage{fontspec}
\usepackage{xeCJK}
\IfFontExistsTF{Times New Roman}
    {\setmainfont{Times New Roman}}
    {\setmainfont{TeX Gyre Termes}}
\IfFontExistsTF{Songti TC}
    {\setCJKmainfont{Songti TC}}
    {\setCJKmainfont{Noto Serif CJK TC}}
\IfFontExistsTF{PingFang TC}
    {\setCJKsansfont{PingFang TC}}
    {\setCJKsansfont{Noto Sans CJK TC}}
% 等寬字用 Latin Modern Mono（只在 TeX Live 樹裡，用檔名讓 kpathsea 找）。
\setmonofont{lmmono10-regular}[
    Extension = .otf,
    AutoFakeBold = 2.5,
]

% ---------------------------------------------------------------
% 本書排版套件（packages/ 目錄）
% ---------------------------------------------------------------
\usepackage{sleek}
\usepackage{sleek-code}
\usepackage{sleek-diagrams}
\usepackage{sleek-recap}
\usepackage{sleek-aki}
\usepackage{sleek-pseudo}   % CLRS 風格的虛擬碼框（algpseudocode）
\usepackage{sleek-math}     % 定義、定理、證明（amsthm）
% 注意：sleek-theorems.sty 不載入，原因見 sleek-math.sty 開頭。
% 注意：sleek-listings.sty 不載入，程式碼一律用 sleek-code.sty 的 minted 環境。

% 章節標題中文化：「第 N 章」
\usepackage{titlesec}
\usepackage{longtable}  % 可以跨頁的長表（附錄 B 的作答提醒總表）
\titleformat{\chapter}{\normalfont\huge\bfseries}{第~\thechapter~章}{1em}{}
\titlespacing*{\chapter}{0pt}{-20pt}{20pt}
% 頁首寫「第 N 章 章名」，不用英文的 CHAPTER。
\renewcommand{\chaptermark}[1]{\markboth{第~\thechapter~章\quad #1}{}}
% 目錄的頁碼欄放寬，四位數頁碼才不會超出右邊界。
\makeatletter
\renewcommand{\@pnumwidth}{2.2em}
% 節號到兩位數（21.10）時，預設的 2.3em 放不下，標題會黏在號碼後面
\renewcommand*\l@section{\@dottedtocline{1}{1.5em}{3.1em}}
\renewcommand*\l@subsection{\@dottedtocline{2}{4.6em}{3.9em}}
\makeatother

% 行內程式碼（\texttt）不能斷行，允許 TeX 多拉開一點字距。
\setlength{\emergencystretch}{3em}

% 詳解：緊接在每一題的 pastexam 框後面
\newenvironment{pastsoln}{\par\smallskip\noindent\textbf{詳解}\quad}{\par\medskip{\color{black!25}\hrule height 0.4pt}\medskip}
\renewcommand{\contentsname}{目錄}
\setlength{\headheight}{14.6pt}
% ⑴⑵…、①②…、㈠㈡… 用中文字型（英文字型沒有這些字）
\xeCJKDeclareCharClass{CJK}{"2460->"24FF, "3220->"3243}
% 填空格：\blank{12} 印成有底線的 (12)
\newcommand{\blank}[1]{\underline{\makebox[4.5em]{(#1)}}}
% 每份考卷的作答說明（附錄「依考卷查題」裡用）
\newenvironment{paperinfo}{\par\small\noindent}{\par\medskip}

\hypersetup{
    pdftitle={資料結構與演算法研究所考古題詳解（三）：他系「資料結構」106–115 學年度},
    pdfauthor={林子祐},
}

\title{資料結構與演算法\\研究所考古題詳解\\[0.3em]{\LARGE （三）他系「資料結構」106--115 學年度}}
\author{林子祐}
\date{\today}

\begin{document}

\maketitle
\romantableofcontents
\chapter*{怎麼用這本書}
\phantomsection\addcontentsline{toc}{chapter}{怎麼用這本書}
\markboth{怎麼用這本書}{}

這一冊收錄資訊類以外系所的研究所入學考試「資料結構」原題：電機、電信、電子、工科海洋、工業工程、生機、電通、資管等，106--115 學年度，每一題都附完整詳解。 這類考卷偏資料結構本身（串列、堆疊、樹、堆積、雜湊、排序、圖的基本演算法），演算法設計和理論的題比較少；資結書書前「學習路徑」的路徑三就是為這類考卷排的讀書順序。 它是
《資料結構與演算法：C 語言、虛擬碼與研究所考試》（以下簡稱\textbf{資結書}）的習題本：
章號和章名跟資結書一樣，第~12~章的題目就是資結書第~12~章雜湊的內容。

\section*{題目怎麼排}
\phantomsection\addcontentsline{toc}{section}{題目怎麼排}

\begin{itemize}
\item \textbf{依主題}：一題歸到它主要考的那一章。 一題同時考兩章時，放在比較吃重的那一章。
\item \textbf{章裡依學校、學年度排}：同一份考卷的題目放在一起，照原卷的題號順序。
\item \textbf{想整份考卷計時練習}：附錄「依考卷查題」列出每一份考卷的每一題在第幾頁。
\item 計算機概論類的考卷只收資料結構與演算法的題目。
\item 組合計數、數學歸納法、數論這類數學題，放在附錄「數學工具」（對應資結書附錄~A）。
\end{itemize}

\section*{題目和詳解}
\phantomsection\addcontentsline{toc}{section}{題目和詳解}

\begin{itemize}
\item 題目照原卷：英文就是英文，記號、配分、選項照錄，錯字也不改。 原卷的圖用 TikZ 重畫。
\item 詳解是本書作者寫的，不是各校公布的答案。 選擇題逐個選項說明對錯，追蹤題寫出每一步，
證明題寫完整的證明，設計題寫虛擬碼、正確性理由和複雜度。 每個數值答案都用程式驗算過。
\item 原題有錯字、條件不足或有兩種讀法時，詳解第一段會說明，並寫出採用的解讀。
\item 詳解裡的「資結書第~N~章」「資結書第~N.M~節」指資結書的章節，虛擬碼的寫法也照資結書
第~1~章的約定。
\end{itemize}

\section*{建議的用法}
\phantomsection\addcontentsline{toc}{section}{建議的用法}

讀完資結書的一章，就做這本書的同一章。 先自己寫，寫完再看詳解；對答案時除了答案本身，也看
詳解的寫法（追蹤畫到哪一步、證明寫到多細），那是閱卷老師看的東西。 考前一兩個月，改用附錄
的「依考卷查題」整份計時做，練時間分配。

\section*{版權}
\phantomsection\addcontentsline{toc}{section}{版權}

依法令舉行之考試試題，依著作權法第~9~條不得為著作權之標的；本書收錄的題目來自各校公開的
研究所入學考試試題。 詳解與排版以 CC BY-NC-SA 4.0 授權。


\input{generated/vol3-body}

\appendix
\titleformat{\chapter}{\normalfont\huge\bfseries}{附錄~\thechapter}{1em}{}
\renewcommand{\chaptermark}[1]{\markboth{附錄~\thechapter\quad #1}{}}
\input{generated/vol3-math}
\input{generated/vol3-papers}

\end{document}
